% Chapter 3

\chapter{Les apports techniques} % Main chapter title

\label{apports} % For referencing the chapter elsewhere, use \ref{apports_techniques}

\lhead{Troisième partie. \emph{Les apports techniques}} % This is for the header on each page - perhaps a shortened title

%----------------------------------------------------------------------------------------
%----------------------------------------------------------------------------------------

Le site Cadremploi.fr est découpé en différentes parties.
En effet, il dispose d'une partie publique, qui est celle dans laquelle un utilisateur peut rechercher une offre en fonction de critères et y postuler.
Une seconde partie, destinée aux recruteurs, leur permet de saisir une offre d'emploi qui sera ensuite consultable depuis la partie client.
J'ai travaillé durant ce stage sur la partie concernant les recruteurs, ou "Espace Recruteur".
L'espace recruteur est aujourd'hui trop ancien, difficilement maintenable et utilisable.
Une refonte totale en était nécessaire.
Mon stage consistait ainsi à participer avec l'équipe Cadremploi à la création d'un nouvel Espace Recruteur, plus moderne.


\section{Cahier des charges et fonctionnement général du cycle de développement}
%----------------------------------------------------------------------------------------
\subsection{Contenu du stage}
\label{sub:Contenu du stage}

Le sujet principal de mon stage concernait la refonte de l'espace recruteur du site cadremploi.fr.
CADREMPLOI.fr a conçu un espace réservé aux professionnels afin de les aider dans leurs campagnes de recrutement.
Grâce à ce service de e-commerce, il leur est possible d'accéder à tous les produits et services qu'offre Cadremploi et de payer directement en ligne avec leur carte de crédit ou par chèque après réception de la facture.
%%%%%%%%%%%TODO%%%%%%%%%%%%%%%%%%%%%%
Cette application web avait déjà fait l'objet d'un développement il y a BLABLA ans et une refonte totale en était nécessaire, de manière à proposer des services nouveaux ainsi qu'un design plus actuel.
Concrètement, les principaux objectifs de ce stage étaient les suivants:
\begin{itemize}
  \item{} Développer de nouvelles fonctionnalités dans un environnement technique dynamique (Play! Scala, ElasticSearch, AngularJS)
  \item{} Maintenir le haut niveau de qualité et de tests présents
  \item{} Participer aux ateliers d'architecture et de cadrages techniques
  \item{} Echanger autour de bonnes pratiques avec les équipes de développement
\end{itemize}
%%%%%%%%%%%TODO%%%%%%%%%%%%%%%%%%%%%%
Les fonctionnalités attendues comprenaient notamment la mise en place du pattern Event Sourcing, que je détaillerai plus tard, de manière à pouvoir suivre de manière rigoureuse l'évolution de chaque utilisateur dans l'espace recruteur, ainsi que la simplification du système de classification des offres en créant un produit \"de base\" auquel se grefferont des options.

%----------------------------------------------------------------------------------------
\subsection{Méthodologie de management}
\label{sub:Méthodologie de management}
\paragraph{}
L'équipe Cadremploi, sur le projet Espace Recruteur notamment, suit une méthodologie Agile.
%%%%%%%%%%%TODO%%%%%%%%%%%%%%%%%%%%%%
% Bref résumé de ce qu'est la méthode Agile

% Rappel, MEP chaque semaine + MEPs exceptionnelles, sprint de 2 semaines
Le système de tickets post-it, des stand-up-meeting était très bien suivi bien que leur pratique ne soit pas un exercice figé.
Le contenu des stand-up-meeting est par exemple passé au cours de mon stage de "Qu'est-ce que j'ai fait?, Qu'est-ce que je vais faire?" à des questions plus pratiques: "Qu'est-ce que j'ai fait de réellement impactant pour le reste de l'équipe? Quels dysfonctionnements ai-je remarqué?".
Finalement, plutôt que l'évolution de chacun dans sa tâche, c'est l'évolution de l'application elle-même qui était alors discutée durant les stands-up, réduisant ainsi le temps de l'exercice, ce qui était nécessaire dans cette équipe de 10 personnes.
\paragraph{}
Les post-its suivaient un cycle de vie dont les étapes rentraient dans les catégories suivantes:
\begin{itemize}
  \item{Nouveau}, il s'agit du statut des tickets qui viennent d'arriver sur le tableau et qui seront donc à traiter prochainement
  \item{En cours}, il s'agit d'un ticket qu'un membre de l'équipe est en train de traiter.
  \item{Recette équipe}, lorsqu'une tâche est terminée, elle doit être validée par un membre de l'équipe; ce système de team review que je traiterai plus tard permet d'éviter de nombreuses erreurs au cours du développement de l'application.
  \item{Recette PO} lorsque le ticket passe l'étape de recette équipe, il arrive en recette PO.
  Le Project Owner s'occupe donc de vérifier que les modifications apportées correspondent bien aux attentes et qu'elles ne perturbent pas les anciennes fonctionnalités.
  Si un problème est rencontré, le PO rédige ses retours et la personne responsable du ticket se charge de corriger les problèmes.
  Le ticket reste dans le même état mais un jeu d'étiquettes "Retours Recette" et "Corrigé" est déposé sur le ticket pour qu'il soit possible de suivre l'évolution des corrections.
  \item{Validé}, le PO valide le ticket et les modifications seront mises en production sous peu
  \item{En attente}, le ticket ne peut être traité pour le moment; une synchronisation avec une autre équipe est nécessaire ou bien des questions restent sans réponse.
\end{itemize}
Ainsi les membres de l'équipe choisissent arbitrairement un ticket qu'il traitent en collant une étiquette les représentant dessus et déplacent le ticket dans la catégorie à laquelle il correspond.
\paragraph{}
Le tableau sur lequel sont affiché les tickets sont plutôt grands et plusieurs méthodes sont utilisés pour permettre à l'équipe de s'y retrouver rapidement.
En effet, en plus des colonnes de catégories dans lesquelles sont rangés les tickets, un code couleur permettait de différencier les tâches techniques des tâches fonctionnelles, ainsi que les tâches traitant de parties de Cadremploi autre que l'espace recruteur (projet de refonte de la Home, ...).
De plus, de petites étiquettes, en plus de celles représentant les membres de l'équipe, sont placées sur les tickets et permettent d'indiquer facilement si un des posts-its représente une tâche qui doit être mise en production dans la semaine (Cible MEP), ou si des retours suite à une recette (équipe ou PO) sont à traiter.
\paragraph{}
Enfin, les post-its représentant nos tâches rédigés sur un tableau se trouvent aussi référencés sur l'outil Redmine, une application web de gestion de projets, qui stocke toutes les informations concernant une tâche donnée.
Ce stockage double de l'information est un peu lourd; en effet, il est nécessaire de déplacer le ticket sur le tableau et d'effectuer le même changement en ligne sur Redmine.
Le point positif est qu'il permet à toute l'équipe, même en télétravail, d'accéder facilement à l'information concernant une tâche, notamment aux différentes discussions ayant eu lieu concernant cette tâche.
Il est cependant nécessaire de conserver l'affichage au tableau puisqu'il est bien plus visuel et permet à l'équipe d'échanger bien plus facilement que si le support était uniquement digital.

\section{Cahier des charges et fonctionnement général du cycle de développement}
%-------------------------------------------------------------------------------
\subsection{Contenu du stage}
\label{sub:Contenu du stage}
Le sujet principal de mon stage concernait ainsi la refonte de l'Espace Recruteur du site cadremploi.fr.
L'application actuelle se présente comme un espace réservé aux professionnels afin de les aider dans leur campagnes de recrutement.
Grâce à ce service de e-commerce, il leur est possible d'accéder à tous les produits et services qu'offre Cadremploi et de payer directement en ligne avec leur carte de crédit ou par chèque après réception de la facture.
Cette application web avait déjà fait l'objet d'un développement il y a plusieurs années et une refonte totale en était nécessaire, de manière à proposer des services nouveaux, un design plus actuel et des éléments de maintenabilité plus importants.
Concrètement, les principaux objectifs de ce stage étaient les suivants:
\begin{itemize}
  \item{} Développer de nouvelles fonctionnalités dans un environnement technique dynamique (Play! Scala, ElasticSearch, AngularJS)
  \item{} Maintenir le haut niveau de qualité et de tests présents
  \item{} Participer aux ateliers d'architecture et de cadrages techniques
  \item{} Echanger autour de bonnes pratiques avec les équipes de développement
\end{itemize}
\begin{figure}[h]
  \begin{center}
    \hspace*{-1in}
    \includegraphics[width=1\textwidth]{Pictures/EspaceRecruteur.png}
  \end{center}
\end{figure}

%-------------------------------------------------------------------------------
\subsection{Méthode de management}
\label{sub:Méthode de management}
%-------------------------------------------------------------------------------
\paragraph{}
L'équipe Cadremploi, sur le projet Espace Recruteur notamment, suit une méthodologie Agile.
Les méthodes de type agile se veulent plus pragmatiques que les méthodes traditionnelles.
Elles impliquent au maximum le client et permettent une grande réactivité à ses demandes.
Les méthodes agiles prônent 4 valeurs fondamentales: l'équipe, qui doit être soudée et dans laquelle la communication est fondamentale, le fonctionnement de l'application, la collaboration avec les clients au delà de la négociation contractuelle ainsi que la flexibilité de la structure du logiciel, acceptant ainsi aisément le changement.
Pour l'équipe Cadremploi, cela se caractérisait par des meetings quotidien devant un tableau rassemblant des tâches demandées par le client écrites sur des post-its, une mise en production des nouveautés développées chaque jeudi, la mise en place de tests unitaires et end-to-end ainsi que des revues de code pour chaque tâche effectuée.
%-------------------------------------------------------------------------------
\paragraph{Le stand-up meeting}
Comme dans de nombreuses méthodes Agiles, la communication au sein de l'équipe passe notamment par un stand-up-meeting quotidien ainsi qu'un affichage mural composé de posts-its représentant les tâches relatives à l'évolution de l'application.
Le système de tickets post-it, des stand-up-meeting était très bien suivi bien que leur pratique ne soit pas un exercice figé.
Le contenu des stand-up-meeting est par exemple passé au cours de mon stage de "Qu'est-ce que j'ai fait?, Qu'est-ce que je vais faire?, Ai-je des problèmes?" à des questions plus pratiques: "Qu'est-ce que j'ai fait de réellement impactant pour le reste de l'équipe? Quels dysfonctionnements ai-je remarqué?".
Finalement, plutôt que l'évolution de chacun dans sa tâche, c'est l'évolution de l'application elle-même qui était alors discutée durant les stands-up, réduisant ainsi le œtemps de l'exercice, ce qui était une nécessité dans cette équipe de 10 personnes.
De plus, suite à un amoncellement de tâches trop important qui produisait un confus global lors de la lecture du tableau, il a été mis en place des règles permettant de limiter le nombre de tickets en cours de traitement, permettant une meilleure organisation de l'équipe.
%-------------------------------------------------------------------------------
\paragraph{Cycle de vie d'une tâche}
Les post-its suivaient un cycle de vie dont les étapes rentraient dans les catégories suivantes:
\begin{itemize}
  \item "Nouveau", il s'agit du statut des tickets qui viennent d'arriver sur le tableau et qui seront donc à traiter prochainement
  \item "En cours", il s'agit d'un ticket qu'un membre de l'équipe est en train de traiter.
  \item "Recette" équipe, lorsqu'une tâche est terminée, elle doit être validée par un membre de l'équipe.
  Ce système de team review permet d'éviter de nombreuses erreurs au cours du développement de l'application: une vérification du code ainsi que des tests écrits sont effectués.
  \item "Recette" PO lorsque le ticket passe l'étape de recette équipe, il arrive en recette PO.
  Le Project Owner s'occupe donc de vérifier que les modifications apportées correspondent bien aux attentes et qu'elles ne perturbent pas les anciennes fonctionnalités.
  Si un problème est rencontré, le PO rédige ses retours et la personne responsable du ticket se charge de corriger les problèmes.
  Le ticket reste dans le même état mais un jeu d'étiquettes "Retours Recette" et "Corrigé" est déposé sur le ticket pour qu'il soit possible de suivre l'évolution des corrections.
  \item "Validé", le PO valide le ticket et les modifications seront mises en production sous peu
  \item "En attente", le ticket ne peut être traité pour le moment; une synchronisation avec une autre équipe est nécessaire ou bien des questions restent sans réponse.
\end{itemize}
Ainsi les membres de l'équipe choisissent arbitrairement un ticket qu'il traitent en collant une étiquette, un autocollant les représentant dessus et déplacent le ticket dans la catégorie à laquelle il correspond.
Les tickets validés sont généralement mis en production chaque semaine, le jeudi, à l'exception de quelques mises en productions (MEP) exceptionnelles.
\paragraph{Organisation du tableau}
Le tableau sur lequel sont affichés les tickets est plutôt grand et plusieurs méthodes sont utilisées pour permettre à l'équipe de s'y retrouver rapidement.
En effet, en plus des colonnes de catégories dans lesquelles sont rangés les tickets, un code couleur permettait de différencier les tâches techniques des tâches fonctionnelles, ainsi que les tâches traitant de parties de Cadremploi autre que l'espace recruteur (projet de refonte de la Home, ...).
De plus, de petites étiquettes, en plus de celles représentant les membres de l'équipe, sont placées sur les tickets et permettent d'indiquer facilement si un des posts-its représente une tâche qui doit être mise en production dans la semaine (Cible MEP), ou si des retours suite à une recette (équipe ou PO) sont à traiter.
\paragraph{Référencement des tâches}
Enfin, les post-its représentant nos tâches rédigées sur un tableau se trouvent aussi référencées sur l'outil Redmine, une application web de gestion de projets, qui stocke toutes les informations concernant une tâche donnée.
Ce stockage double de l'information est un peu lourd; en effet, il est nécessaire de déplacer le ticket sur le tableau et d'effectuer le même changement en ligne sur Redmine.
Le point positif est qu'il permet à toute l'équipe, même en télétravail, d'accéder facilement à l'information concernant une tâche, notamment aux différentes discussions ayant eues lieu concernant cette tâche.
Il est cependant nécessaire de conserver l'affichage au tableau puisqu'il est bien plus visuel et permet à l'équipe d'échanger bien plus facilement que si le support était uniquement digital.

\section{Aperçu des technologies et des techniques utilisées}

\subsection{Langages}
\label{sub:Langages}
\subsubsection{Backend}
\label{subs:Backend}
\begin{figure}[h]
  \begin{center}
    \includegraphics[width=0.3\textwidth]{Pictures/play_logo.png}
    \hspace{1in}
    \includegraphics[width=0.15\textwidth]{Pictures/scala_logo.png}
  \end{center}
\end{figure}

Le nouvel espace recruteur est géré en backend via le framework Play! dans sa dernière version (2.4.2) et sous le langage Scala.
Il s'agit d'un couple qui permet d'écrire du code rapidement et de façon maintenable.
Play est l'un des framework les plus connus dans le monde Java et offre de nombreux avantages:
\begin{itemize}
  \item Il supporte le rechargement à chaud.
  Il suffit de lancer Play via sa console en mode développement et il prendra automatiquement en compte à chaud les changements effectués sur le code, mais aussi les templates ou le routage.
  Cela contribue grandement au gain de productivité qu'offre Play.
  \item En plus de s’appuyer sur du code à typage statique, Play propose la sécurité du typage à d’autres endroits et notamment sur les templates ou sur le routage des différents contrôleurs.
  Un certain nombre de problèmes sont alors mis en lumière directement à l’étape de la compilation.
  \item Il permet d'exécuter des tests, notamment sur la couche web à plusieurs niveaux:
  On peut par exemple tester un contrôleur en démarrant un serveur web (donc via HTTP) ou, sans le démarrer, en appelant simplement le contrôleur avec le contexte qui va bien, le tout de manière simple et rapide à l’exécution.
  \item Il est sans état et basé sur des entrées sorties non bloquantes et permet ainsi une capacité à monter en charge très intéressante
  \item Il supporte nativement REST, JSON, Websocket entre autres et se présente donc comme un framework moderne
\end{itemize}
Un des points négatifs de ce choix en backend est l'utilisation de fait obligatoire de SBT qui nous a posé de nombreux problèmes.
En effet, cet outil était souvent lent, compliqué d'utilisation et ralentissait souvent le développement notamment via ses résolutions de dépendances interminables.
\subsubsection{Frontend}
\label{subs:Frontend}
\begin{figure}[h]
  \begin{center}
    \includegraphics[width=0.4\textwidth]{Pictures/angular_logo.png}
  \end{center}
\end{figure}
%% TODO
Du côté front, l'équipe a utilisé le framework AngularJS ainsi que la librairie D3.js qui permettent des animations dynamiques.
Ce choix de technologie peut être soutenu pour les points suivants:
\begin{itemize}
  \item Angular permet une maintenabilité de l'application aisée, notamment en demandant une structure de type MVC au développeur, ainsi qu'en utilisant du HTMl, qui est déclaratif, pour définir l'interface utilisateur.
  \item Il s'agit de plus d'un framework flexible puisqu'il est possible de composer une application en
\end{itemize}

\subsection{Environnements et outils}
\paragraph{IntelliJ}
\label{par:IntelliJ}
\begin{wrapfigure}{l}{0.15\textwidth}
  \vspace{-2.5em}
  \begin{center}
    \includegraphics[width=0.15\textwidth]{Pictures/intellij_logo.png}
  \end{center}
\end{wrapfigure}
Le développement de l'espace recruteur s'est fait sous le framework Play! et IntelliJ Idea était l'IDE utilisé par l'équipe pour développer, puisqu'il offre une bonne intégration de ce framework.
J'avais déjà souvent utilisé cet IDE pour différents projets que j'ai eu à rendre tout au long de ma scolarité à l'EISTI, ainsi son utilisation ne m'a posé aucun problème.
\paragraph{Git}
\label{par:Git}
L'équipe étant nombreuse et l'application étant sujete à grossir avec le temps, un outil de gestion de versions a été utilisée.
Git est l'outil le plus répandu et le plus efficace connu, c'est ce que l'équipe Cadremploi a utilisé.
Il s'agissait aussi d'un outil que j'avais beaucoup utilisé mais de manière très simpliste et j'ai appris à m'en servir d'une toute autre façon, bien plus complète, pendant ce stage.
\paragraph{Jenkins}
\label{par:Jenkins}
\begin{wrapfigure}{rH}{0.25\textwidth}
  \begin{center}
    \includegraphics[width=0.15\textwidth]{Pictures/jenkins_logo.png}
  \end{center}
\end{wrapfigure}
L'équipe Cadremploi utilise aussi Jenkins comme outil d'intégration continue pour déployer les nouveautés sur les différents environnement de production utilisés (recette, intégration), les mises en production et préproductions n'étant pas gérées directement par l'équipe.
Cet outil avait été brièvement présenté lors du cours d'Outils de Développement mais jamais utilisé.
Jenkins est un outil open source permettant de compiler et de tester un projet de manière continue.
Il aide de fait les développeurs à intégrer facilement des changements à une application.
Bien que nous avions été tenté de mettre en place un Jenkins avec mon équipe de Projet de Fin d'Étude, cela ne s'est jamais fait par manque de temps et puisqu'un besoin concret ne s'est jamais montré.
Je n'ai pas eu de mal à l'utiliser puisqu'une fois configuré, cet outil est réelement simple d'utilisation.
Plusieurs environnements de développement sont utilisés par Cadremploi comme pour les autres équipes de Figaro Classifieds:
\begin{itemize}
  \item{Environnement de développement}, où les développeurs postent régulièrement les améliorations qu'ils mettent en place. Jenkins n'est pas utilisé pour cet environnement puisque les changements y sont publiés directement via Git.
  \item{Environnement de recette}, où la Project Owner vérifie que les changements apportés fonctionnent effectivement et les valide.
  \item{Environnement de pré-production} qui est un environnement qui ressemble autant que possible à l'environnement de production et où sont surtout testés les architectures
  \item{Environnement de production} qui est la dernière étape, contenant l'application utlisée réelement par les utilisateurs
  \item{Environnement de'intégration}, il s'agit d'un environnement supplémentaire qui est préféré pour les tests fait sur l'infrastructure du projet.
\end{itemize}
J'ai été habitué durant ma scolarité à l'EISTI à utiliser la majorité des outils présentés précédemment puisque j'étais uniquement étranger à Jenkins.
J'ai malgré tout grandement progressé dans mon utilisation de ces outils et particulièrement pour Git que j'ai redécouvert; le projet auquel je participais differais grandement en taille de ceux auquels j'avais contribué auparavant.
%Potentiellement, Rundeck, Kibana, Puppet.

\section{Une équipe nombreuse}
\label{sec:Une équipe nombreuse}
Cadremploi, a mon arrivée, comptait alors 10 membres dont 9 développeurs.
Il s'agit d'une équipe d'une taille importante pour un projet Agile.
Ainsi, de l'organisation ainsi qu'une bonne coordination étaient nécessaires.

\subsection{Revue de code}
\label{sub:Revue de code}
Le développement d'applications demande très souvent à une équipe de personnes de travailler ensemble.
Plus l'équipe grandit plus il est compliqué de garantir une coordination correcte et une bonne organisation.
Les équipes d'ingénieur sont spécialement soumises à ce genre de problèmes puisque du code est quotidiennement partagé entre plusieurs personnes au sein d'une entreprise.
La revue de code aide à partager le savoir et les bonnes pratiques.
\paragraph{}
Le flux de production de l'équipe Cadremploi demande à chaque membre de l'équipe à passer un ticket terminé dans un état "R7 Equipe".
Le ticket doit alors être revu par un autre développeur de l'équipe.
Ce dernier doit analyser les changements fait au code, s'assurer qu'ils n'introduisent pas de troubles dans l'application et vérifier que le code produit suit bien les normes de code imposées par l'équipe.
\paragraph{}
La revue de code par d'autres membres de l'équipe offre de nombreux avantages autres qu'une simple vérification du code.
En effet, il est clair que n'avoir qu'un unique membre de l'équipe responsable d'un point critique sur l'application est dangereux.
La revue de code distribue le savoir au sein de l'équipe et permet d'informer plusieurs membres d'un changement.
De plus, elle stimule les conversations portant sur la structure du code, l'architecture de l'application et permet ainsi à de nouveaux venus, comme je l'étais, de s'adapter rapidement et donc de devenir productif plus rapidement.
\paragraph{}
Il ne m'est pas possible de réelement comparer ce mode de fonctionnement avec mes projets précédents.
En effet, mes projets passés se faisaient dans des équipes plus réduites, et le temps manquait pour la revue de code.
Cela ne nous empêchait pas de suivre la règle du boyscoot, qui prescrit de toujours laisser une fonction, un module dans un meilleur état que celui dans lequel on l'a trouvé, et qui évite la "possession" du code par un membre de l'équipe.
Cette règle est plus ou moins comprise dans la revue de code de Cadremploi qui offre plus d'avantages encore.

\subsection{Utilisation de Git}
\label{sub:utilisation de Git}
Git est un outil de contrôle de version flexible et puissant.
Il propose énormément d'options de workflow et il est parfois compliqué
Git is a flexible and powerful version control system. While Git offers significant functionality over legacy centralized tools like CVS and Subversion, it also presents so many options for workflow that it can be difficult to determine what is the best method to commit code to a project. The following are the guidelines I like to use for most software projects contained within a Git repository. They aren't applicable to every Git project (especially those hosted on drupal.org or GitHub), but I've found that they help ensure that our own projects end up with a reasonable repository history.

\section{CQRS: Command Query Responsibility Seggregation}
\label{sec:CQRS: Command Query Responsibility Seggregation}
Le pattern CQRS repose sur le principe de séparation entre l'écriture et la lecture de l'information.
C'est un pattern que nous avons suivi lors du développement de l'espace recruteur puisque la séparations des composants de traitement (les "commands") et de restitution (les "queries") de l'information offrait une architecture très intéressante de laquelle nous avons tiré de nombreux bénéfices tels que la suppression du risque d'effets de bord ou l'allègement des classes de service.
\subsection{Le pattern Command}
\label{sub:Le pattern Command}
Une commande représente une action destinée à être exécutée.
Dans le cas de l'espace recruteur du site Cadremploi.fr, ces commandes ont pour but de stocker en base les événements altérant le système.
Concrètement, ce pattern consiste à englober un appel de méthode dans un objet de type Command qui appelle un objet de type CommandHandler qui appelera lui la-dite méthode.
Bien que cela augmente les lignes de code (il faut grossièrement 2 classes pour appeler une méthode), le code que propose le pattern Command est plus expressif, mieux découpé et simple à tester.
\subsection{Le pattern Query}
\label{sub:Le pattern Query}
La partie Query de ce modèle se base sur le fait que les objets du modèle sont volumineux et qu'il est possible de s'en passer.
En effet le besoin représenté ici est celui d'une lecture dans un cas d'utilisation bien précis.
Le pattern Query consiste alors à exécuter une requête précise en base et de restituer un objet DTO concis qui pourra être utilisé directement.
Le modèle CQRS de base peut s'avérer problématique puisque la grosse majorité des sollicitations du système se fait sur cette partie.
% Pour le site Cadremploi, cette requête est une requête ElasticSearch et est donc extrêmement rapide.
% La partie de Read de l'application se résume ainsi à la simple exécution d'une requête ES et au remplissage de DTO.

\section{Event sourcing}
% http://martinfowler.com/eaaDev/EventSourcing.html
% https://github.com/eventstore/eventstore/wiki/Event-Sourcing-Basics
Une des demandes pour le nouvel espace recruteur est de pouvoir suivre avec précision les actions de l'utilisateur.
Ainsi chaque modification qu'il apportera à une de ses offres ou à son profil par exemple doit être enregistrée.
%----------------------------------------------------------------------------------------
\textit{"L'idée fondamentale derrière l'Event Sourcing est d'assurer que chaque modification de l'état d'une application est capturée dans un objet événement et que ces événements sont eux-mêmes stockés dans l'ordre dans lequel ils ont été appliqués à l'application."} (Martin Fowler)
%----------------------------------------------------------------------------------------
\subsection{Fonctionnement désiré}
\label{sub:Fonctionnement desire}
Il est nécessaire de pouvoir suivre l'évolution de l'annonce que créé un utilisateur.
Concrètement, à chaque fois que l'utilisateur modifie son annonce, on veut générer un événement qui informe notre système de cette modification.
Toutes ces données pourront ainsi être utilisées pour faire des statistiques ainsi que des contrôles.
Une des utilités de ce suivi est par exemple de pouvoir contrôler les éditions abusives d'annonces déjà en ligne.
Il est ainsi possible de suivre de très près la façon dont l'utilisateur procède, ou les champs qui posent problème.

%----------------------------------------------------------------------------------------
\subsection{Fonctionnement technique}
%----------------------------------------------------------------------------------------
\subsubsection{Enregistrement de l'événement}
\label{subs:Enregistrement de l'evenement}
Concrètement, derrière chaque champ formulaire que l'utilisateur remplit, un watcher d'Angular est présent.
Ainsi lorsque ce watcher indique que le champ a été modifié, on récupère la valeur qu'il contient et envoyons cette demande de modification vers le backend.
Cette "demande de modification" est traitée comme une commande et génère alors un ou plusieurs événements qui résument la modification effectuée et sont stockés dans une base de donnée, conformément au modèle de Write du pattern CQRS décrit ci-dessus.
Chacune des modifications faite par l'utilisateur est ainsi enregistrée, et il est possible de connaître l'état du système à tout instant t donné en interrogeant la base de données.
%----------------------------------------------------------------------------------------
\subsubsection{Utilisation de l'événement}
\label{subs:Utilisation de l'evenement}
Si l'enregistrement en base se passe bien, l'événement est ensuite publié sur un EventBus, un objet Akka qui permet entre autre d'envoyer un message à des groupes d'acteurs, pour être écouté par des objets de type Projection.
Ces projections écoutent seulement une partie des événements selon leurs responsabilité et modifient en conséquence un objet DTO qu'elles indexent.
Ce traitement est fait de manière asynchrone, via le système d'acteurs Akka.
Il permet donc d'optimiser le modèle CQRS vu précédemment: la commande est exécutée rapidement puisque l'événement est sauvegardé de suite et le reste du traitement est effectué de manière asynchrone.
Cela permet d'informer quasi instantannément l'utilisateur que son action a été prise en compte.
D'autre part, en acceptant de rendre asynchrone et donc pas forcément instantannée la mise à jour de la partie Read du modèle, on améliore grandement la performance ressentie par l'utilisateur.
Le système est donc tenu à jour en lecture et ainsi, à l'appel d'une query, les projections peuvent être interrogées pour récupérer la donnée de manière très rapide.


\section{Domain Driven Development}


%----------------------------------------------------------------------------------------
%----------------------------------------------------------------------------------------
\section{Organisation}
comparaison entre les plannings prévisionnels et réels

%----------------------------------------------------------------------------------------

%----------------------------------------------------------------------------------------
%----------------------------------------------------------------------------------------

\section{Auto-évaluation}
respect des délais, autonomie, qualité du travail, apports à l'équipe

%----------------------------------------------------------------------------------------
%----------------------------------------------------------------------------------------

\section{Résultats et prolongements possibles}
le site tourne et continuera de fonctionner encore

%----------------------------------------------------------------------------------------
